\documentclass[12pt,letterpaper]{article}

\usepackage[utf8]{inputenc}
\usepackage[spanish]{babel}
\usepackage[margin=2.5cm]{geometry}
\usepackage{amsmath}
\usepackage{amssymb}
\usepackage{graphicx}
\usepackage{tikz}
\usetikzlibrary{shapes.geometric, arrows.meta, positioning, calc}

\title{Multiplicadores de Lagrange aplicados a la Teoría del Consumidor\\
\large Optimización de presupuesto y diagramas de flujo del programa en C++}
\author{Angel Manuel Ruvalcaba Garcia}
\date{\today}

% Estilos para los diagramas de flujo
\tikzstyle{inicio} = [ellipse, draw, fill=blue!15, minimum width=2.6cm, minimum height=1cm, align=center]
\tikzstyle{proceso} = [rectangle, draw, fill=orange!15, minimum width=4.6cm, minimum height=1cm, align=center, text width=4.4cm]
\tikzstyle{decision} = [diamond, draw, fill=yellow!20, minimum width=2.6cm, minimum height=1.4cm, align=center, text width=2.4cm, inner sep=1pt]
\tikzstyle{entradasalida} = [trapezium, draw, fill=green!15, trapezium left angle=70, trapezium right angle=110, minimum width=3cm, minimum height=1cm, align=center, text width=3cm]
\tikzstyle{flecha} = [thick, -{Stealth[length=2.5mm]}]
\tikzstyle{procesoFormula} = [rectangle, draw, fill=orange!15, minimum width=4.6cm, minimum height=1.7cm, align=center, text width=4.4cm]

\begin{document}
\maketitle

\section{¿Qué es un multiplicador de Lagrange?}

Un multiplicador de Lagrange ($\lambda$) es una variable auxiliar que se introduce para resolver problemas de \textbf{optimización con restricciones}: maximizar (o minimizar) una función sujeta a que las variables cumplan una condición adicional (una igualdad).

En vez de resolver el problema restringido directamente, se construye una nueva función --el \textbf{Lagrangiano}-- que combina la función objetivo y la restricción:
\[
L(x_1,\dots,x_n,\lambda) = f(x_1,\dots,x_n) + \lambda \big(c - g(x_1,\dots,x_n)\big)
\]
donde $f$ es la función que se quiere optimizar y $g(x_1,\dots,x_n) = c$ es la restricción. El multiplicador $\lambda$ mide la \textbf{sensibilidad del óptimo ante un cambio marginal en la restricción}: cuánto mejora (o empeora) el valor óptimo de $f$ si la restricción se relaja o se aprieta ligeramente.

\subsection{¿Para qué se utiliza?}

Se utiliza siempre que existan dos elementos: algo que se quiere maximizar o minimizar, y un límite (recurso escaso) que no se puede exceder. Ejemplos típicos:
\begin{itemize}
    \item Un consumidor que maximiza su utilidad sujeto a un presupuesto limitado (nuestro caso).
    \item Una empresa que minimiza costos de producción sujeta a producir cierta cantidad.
    \item Un ingeniero que maximiza la resistencia de una estructura sujeta a un límite de material.
\end{itemize}

El método convierte un problema restringido en un sistema de ecuaciones (derivadas parciales igualadas a cero) que se puede resolver con álgebra estándar.

\section{Aplicación a la teoría del consumidor}

\subsection{Planteamiento del problema}

Un consumidor quiere maximizar su utilidad $U(x_1,\dots,x_n)$ al consumir $n$ bienes, sujeto a no exceder su ingreso disponible $I$:
\[
\text{Maximizar: } U(x_1,\dots,x_n) \qquad
\text{Sujeto a: } \sum_{i=1}^{n} p_i x_i = I
\]
donde $x_i$ es la cantidad consumida del bien $i$, y $p_i$ su precio.

\subsection{El Lagrangiano}

\[
L(x_1,\dots,x_n,\lambda) = U(x_1,\dots,x_n) + \lambda\left(I - \sum_{i=1}^n p_i x_i\right)
\]

\subsection{Condiciones de primer orden}

\[
\frac{\partial L}{\partial x_i} = \frac{\partial U}{\partial x_i} - \lambda p_i = 0 \quad \text{para cada } i
\qquad
\frac{\partial L}{\partial \lambda} = I - \sum_i p_i x_i = 0
\]

De aquí se obtiene la condición de óptimo del consumidor:
\[
\frac{\partial U / \partial x_i}{p_i} = \lambda \quad \text{para todo } i
\]
Es decir, en el óptimo la utilidad marginal por peso gastado debe ser igual en todos los bienes. $\lambda$ representa la utilidad marginal del ingreso: cuánta utilidad adicional se obtendría con un peso más de presupuesto.

\subsection{Caso Cobb-Douglas: solución cerrada}

Para que el programa en C++ pueda calcular el óptimo sin un método numérico iterativo, se modela la utilidad como Cobb-Douglas generalizada:
\[
U(x_1,\dots,x_n) = \prod_{i=1}^{n} x_i^{a_i}, \qquad a_i > 0
\]
donde $a_i$ representa el peso o preferencia relativa del bien $i$. Resolviendo el sistema de condiciones de primer orden se llega a la fórmula cerrada que implementa el programa:
\[
\boxed{x_i^{*} = \dfrac{a_i}{\sum_{j=1}^n a_j} \cdot \dfrac{I}{p_i}}
\]
Esto significa que el gasto óptimo en cada bien es una fracción fija del ingreso, proporcional a su preferencia relativa. El multiplicador se recupera con:
\[
\lambda = \frac{a_i \cdot U^{*}}{x_i^{*} \cdot p_i}
\]
usando cualquier bien $i$ (el resultado debe ser el mismo, lo que sirve como verificación del cálculo).

\section{Estructuras de datos del programa}

\begin{verbatim}
struct Bien {
    string nombre;
    double precio;         // p_i
    double preferencia;    // a_i
    double cantidadOptima; // x_i*
    double gastoOptimo;    // p_i * x_i*
};

class ProblemaPresupuesto {
    vector<Bien> bienes;
    double ingreso;           // I
    double sumaPreferencias;  // suma de a_j
    double lambda;
    double utilidadOptima;

    void agregarBien(nombre, precio, preferencia);
    void establecerIngreso(I);
    void resolver();
    double calcularUtilidad();
    double calcularLambda();
    bool validarRestriccion();
    void mostrarResultados();
};
\end{verbatim}

\section{Diagrama de flujo: \texttt{main()}}

\begin{center}
\begin{tikzpicture}[font=\footnotesize]
\node (inicio)   [inicio]        at (0,0)      {Inicio};
\node (crear)    [proceso]       at (0,-1.4)   {Crear objeto\\\texttt{ProblemaPresupuesto}};
\node (leern)    [entradasalida] at (0,-2.8)   {Leer $n$\\(número de bienes)};
\node (decision) [decision]      at (0,-7.2)   {¿Quedan\\bienes por\\leer?};
\node (leerbien) [entradasalida] at (6,-7.2)   {Leer nombre,\\precio $p_i$,\\preferencia $a_i$};
\node (agregar)  [proceso]       at (6,-8.8)   {\texttt{agregarBien()}};
\node (leerI)    [entradasalida] at (0,-11.6)  {Leer ingreso $I$};
\node (estI)     [proceso]       at (0,-13.0)  {\texttt{establecerIngreso()}};
\node (resolver) [proceso]       at (0,-14.4)  {\texttt{resolver()}};
\node (validar)  [proceso]       at (0,-15.8)  {\texttt{validarRestriccion()}};
\node (mostrar)  [proceso]       at (0,-17.2)  {\texttt{mostrarResultados()}};
\node (fin)      [inicio]        at (0,-18.6)  {Fin};

\draw [flecha] (inicio) -- (crear);
\draw [flecha] (crear) -- (leern);
\draw [flecha] (leern) -- (decision);
\draw [flecha] (decision) -- node[above] {Sí} (leerbien);
\draw [flecha] (leerbien) -- (agregar);
\draw [flecha] (agregar.south) -- (6,-19.8) -- (-2.5,-19.8) -- (-2.5,-7.2) -- (decision.west);
\draw [flecha] (decision) -- node[left] {No} (leerI);
\draw [flecha] (leerI) -- (estI);
\draw [flecha] (estI) -- (resolver);
\draw [flecha] (resolver) -- (validar);
\draw [flecha] (validar) -- (mostrar);
\draw [flecha] (mostrar) -- (fin);
\end{tikzpicture}
\end{center}

\section{Diagrama de flujo: \texttt{resolver()}}

\begin{center}
\begin{tikzpicture}[font=\footnotesize]
\node (inicio2) [inicio]         at (0,0)     {Inicio\\\texttt{resolver()}};
\node (suma)    [proceso]        at (0,-1.4)  {\texttt{sumaPreferencias}\\$= \sum_j a_j$};
\node (loop)    [decision]       at (0,-5.8)  {¿Quedan\\bienes por\\procesar?};
\node (calcx)   [procesoFormula] at (6.2,-5.8) {$x_i^{*} = \dfrac{a_i}{\sum_j a_j} \cdot \dfrac{I}{p_i}$};
\node (calcg)   [procesoFormula] at (6.2,-8.0) {$\text{gastoOptimo}_i = p_i \cdot x_i^{*}$};
\node (util)    [proceso]        at (0,-10.2) {\texttt{calcularUtilidad()}\\$U^{*} = \prod_i (x_i^{*})^{a_i}$};
\node (lam)     [procesoFormula] at (0,-12.0) {\texttt{calcularLambda()}\\$\lambda = \dfrac{a_i \cdot U^{*}}{x_i^{*}\cdot p_i}$};
\node (fin2)    [inicio]         at (0,-13.8) {Fin\\\texttt{resolver()}};

\draw [flecha] (inicio2) -- (suma);
\draw [flecha] (suma) -- (loop);
\draw [flecha] (loop) -- node[above] {Sí} (calcx);
\draw [flecha] (calcx) -- (calcg);
\draw [flecha] (calcg.south) -- (6.2,-15.0) -- (-2.5,-15.0) -- (-2.5,-5.8) -- (loop.west);
\draw [flecha] (loop) -- node[left] {No} (util);
\draw [flecha] (util) -- (lam);
\draw [flecha] (lam) -- (fin2);
\end{tikzpicture}
\end{center}

\end{document}
